% !TeX root = ../prompt-engineering-guide-for-beginners.tex

%%%%%%%%%%%%%%%%%%%%%%%%%%%%%%%%%%
%% 第七章：提示工程进阶篇
%%%%%%%%%%%%%%%%%%%%%%%%%%%%%%%%%%
\chapter{提示工程进阶：让回答更靠谱的增强技巧}
\label{ch:prompting-intermediate}

上一章介绍的基础四板斧，已经能应对你在办公和生活中的多数日常任务。
但如果你遇到需要极其严谨或关乎重要决策的任务，并且想让 AI 的回答彻底摆脱偶尔的“胡说八道”，就需要用到下面这六种进阶增强技巧。
它们是基础技巧的「升级补丁」，能大幅提升模型的可靠性。

\section{自我一致性（Self-Consistency）}
\label{sec:self-consistency}

上一章我们学到了链式思考（CoT，参见第~\ref{sec:cot}~节），通过让 AI 展示推理过程来提高准确率。但即便如此，一条孤立的思维链依然有可能「想偏」，一旦前面的逻辑出现谬误，后面的推理就会全线崩溃。

自我一致性就是用来对冲这层风险的：它的做法是让 AI 针对同一个问题，在底层通过调整 Temperature（温度树）或者 Top-p 参数，强行走出多条不同的推导计算分支（解码路径），然后利用集成学习式的“多数投票（Majority Vote）”机制，对得出的多个最终大纲进行比对，采纳绝大多数路径都指向的那一个答案。

\textit{类比：同一道数学题让实习生用三种方法做，三种方法都得到同一个答案，你就更放心了。}

\textbf{使用方法：}

\begin{promptbox}
以下是一个逻辑推理题：

一家公司有 100 名员工，其中 60\% 会骑自行车，70\% 会游泳，80\% 会打篮球。请问至少有多少员工同时会这三项技能？

请用三种不同的推理方法来解答这道题，然后对比三种方法的结果，给出最终答案。
\end{promptbox}

当你面对重要决策或复杂推理题时，使用自我一致性可以显著降低 AI 出错的概率。

\section{生成知识提示（Generated Knowledge Prompting）}

面对相对生僻或专业度极高的问题，直接发问很容易迫使 AI 开始“盲答”或捏造。此时，更理智的策略是：先命令 AI 去它庞大的底层数据库中，将相关的专业知识点系统性地提取出来，完成一次「闭卷默写」，然后再让它拿着这份自己默写出来的知识库去推演解答。

\textit{类比：考试前让实习生先把知识点默写一遍，再开始答题，正确率更高。}

这个技巧分两步：

\textbf{第一步：生成知识}
\begin{promptbox}
请列出关于「地中海饮食」的 5 个核心知识点，包括其定义、主要食物组成、健康益处、适用人群和常见误区。
\end{promptbox}

\textbf{第二步：基于知识回答问题}
\begin{promptbox}
基于以上知识点，请帮一位 45 岁、有轻度高血压的上班族制定一份为期一周的地中海饮食方案。
\end{promptbox}

通过两步走，AI 的回答会更加有理有据，而不是凭空编造。

当然，你也可以把两步合并到一条提示词里：

\begin{promptbox}
请先回顾关于「地中海饮食」的核心知识，然后基于这些知识，为一位 45 岁有轻度高血压的上班族制定一周饮食方案。
\end{promptbox}

\section{提示链（Prompt Chaining）}

有些任务太复杂，一条提示词搞不定。提示链的思路是：把一个大任务拆成多个小步骤，每一步是一条独立的提示词，前一步的输出直接作为下一步的输入。

\textit{类比：不是让实习生一口气写完一份商业计划书，而是先做市场调研，再写竞品分析，然后定策略，最后组装成文。}

\textbf{实际案例：写一篇产品评测文章}

\begin{promptbox}
\textbf{第一步}：「请帮我列出 iPhone 16 的 10 个核心特性，按照硬件、软件、设计三个维度分类。」

\textbf{第二步}：「基于上面的特性列表，请从一个普通消费者的角度，对每个特性给出优点和不足的评价。」

\textbf{第三步}：「请把上面的分析整合成一篇 1500 字的产品评测文章，结构为：开头引言、核心亮点、值得吐槽的地方、购买建议、总结。语气客观但不无聊。」
\end{promptbox}

提示链的优势在于：

\begin{itemize}
  \item 每一步的任务都足够简单，AI 更容易做好
  \item 你可以在每一步结束后检查结果，发现问题及时纠正
  \item 整体输出的质量通常比一步到位要高很多
\end{itemize}

\section{方向性刺激提示（Directional Stimulus Prompting）}

有时候你手头有一份繁杂的初始材料，你既不想直接把答案喂到 AI 嘴里限制它的发挥，又担心它的思路宛如脱缰野马。这时候，方向性刺激提示就能派上用场：在提示词的结尾埋入一个「高亮线索」或「思考引子」，牵引 AI 的视觉重心。

\textit{类比：侦探导师不直接告诉学徒凶手是谁，而是指了指带血的门把手说「去查查这个把手的指纹」，学徒就能顺藤摸瓜。}

\textbf{不加暗示：}
\begin{promptbox}
请分析这家公司最近业绩下滑的原因。

（附上公司的财务报告）
\end{promptbox}

\textbf{加了暗示：}
\begin{promptbox}
请分析这家公司最近业绩下滑的原因。重点关注研发支出的变化趋势和新品上市节奏。

（附上公司的财务报告）
\end{promptbox}

通过给出一些方向性的提示，AI 就不会在无关紧要的内容上浪费篇幅，而是聚焦在你真正关心的问题上。

\section{Active-Prompt}

Active-Prompt 的核心想法是：先让 AI 回答一批问题，找出它回答得不一致或者不确定的那些（这些就是它的「薄弱环节」），然后针对这些难题，人工补充示例和注释，再让 AI 重新回答。

\textit{类比：先让实习生做一套模拟题，看他哪里老错，然后专门针对错题讲解。}

这个技巧的操作流程：

\begin{enumerate}
  \item 给 AI 一批问题，让它回答
  \item 检查哪些回答不够好或者前后不一致
  \item 针对这些问题，提供详细的示例和解释
  \item 把补充后的提示词再给 AI，让它重新回答
\end{enumerate}

来看一个实际场景。假设你需要 AI 帮你做客户反馈的情感分类。你先让它分类一批反馈，发现大部分「正面」和「负面」的分类都很准确，但碰到带有反讽语气的评论（比如「哦，你们的售后真是太棒了呢」），AI 总是判断错误。

这时候，你就可以针对这类「反讽」场景，补充几个带标注的示例：

\begin{promptbox}
以下是一些客户反馈的情感分类示例，请特别注意带有反讽或讽刺语气的表述：

反馈：「产品质量非常好，用了一年都没问题。」
分类：正面

反馈：「哦，你们的售后真是太棒了呢，打了三天电话都没人接。」
分类：负面（反讽）

反馈：「谢谢你们的快递，三周才到，我还以为包裹去环游世界了。」
分类：负面（反讽）

现在请对以下反馈进行情感分类：
（粘贴新的反馈内容）
\end{promptbox}

通过这种「先诊断弱点，再针对性补强」的方式，你可以用最小的工作量获得最大的质量提升。这个方法特别适合需要批量处理类似任务的场景，比如数据标注、分类文档，或者批量生成特定格式的内容。

\section{检索增强生成（RAG）}
\label{sec:rag}

检索增强生成（Retrieval-Augmented Generation，简称 RAG）是目前非常热门的一个技术。它的核心思路是：让 AI 在回答问题之前，先从你提供的文档或数据库里检索相关内容，然后基于检索到的内容来回答。

\textit{类比：不让实习生凭记忆回答客户问题，而是要求他先查公司知识库，找到依据再回答。}

为什么需要 RAG？因为 AI 的训练数据有截止日期，它不知道最新的信息；它也不了解你公司内部的规章制度、产品文档等私有知识。通过 RAG，你可以让 AI 「看到」这些内容，大幅减少它胡说八道的概率。

\textbf{简单的 RAG（手动版）：}

你可以直接在提示词里附上参考资料：

\begin{promptbox}
请根据以下产品文档，回答客户的问题。如果文档里没有提到的信息，请直接说「这个问题我需要确认后再回复您」。

【产品文档内容】

（粘贴你的文档内容）

客户问题：你们的退货政策是什么？
\end{promptbox}

\textbf{自动化的 RAG：}

更高级的大型 RAG 工程系统：利用 Embedding 模型把你的文档切成数百个段落碎块（Chunks）并转化为高维向量数组存入向量数据库（Vector DB），当用户提问时，系统会先启动距离相似度计算（如余弦相似度）自动检索最匹配的 TOP 5 文档块，再把它们打包当成背景 Context 发射给大模型融合宣读。市面上热门的各类「企业级专属知识库」及各种智能问答机器人的核心底层基石清一色都是 RAG。

对于绝大多数没有代码基础的零基础用户而言，这种「手动粘贴附带文档资料」的新手版 RAG，就已经能在工作流中产生降维打击般的威力了。你需要永远信奉一个原则：\textbf{与其让模型全凭它的赛博记忆去解答问题，不如直接把相关说明书「喂」进它的上下文里，这也是避免 AI 职场翻车的最强保险丝。}

这六种进阶技巧的共同目标都是一样的：通过更精准的引导方式，让 AI 给出确定性更强、更可靠的回答。不需要一次性全部记住，先留个印象，等在实践中遇到了瓶颈，再回到这一章查阅对应的技巧即可。如果你好奇提示工程还能如何进一步拓展，下一章将介绍一些更前沿的高阶方法。
